\documentclass[11pt,a4paper]{article}

\usepackage[a4paper,top=18mm,bottom=19mm,right=19mm,left=19mm]{geometry}
\usepackage{amsmath,amssymb}
\usepackage{xcolor}
\usepackage{array}
\usepackage{enumitem}
\usepackage[explicit]{titlesec}
\usepackage[most]{tcolorbox}
\usepackage{fancyhdr}
\usepackage[colorlinks=true,linkcolor=blue!45!black]{hyperref}
\usepackage{xepersian}

\settextfont[
  Path={../booklet/font/},
  UprightFont={IRLotusICEE.ttf},
  BoldFont={IRLotusICEE_Bold.ttf},
  BoldItalicFont={IRLotusICEE_BoldIranic.ttf},
  ItalicFont={IRLotusICEE_Iranic.ttf},
  FontFace={b}{n}{IRLotusICEE_Bold.ttf},
  FontFace={bx}{n}{IRLotusICEE_Bold.ttf},
  FontFace={b}{it}{IRLotusICEE_BoldIranic.ttf},
  FontFace={bx}{it}{IRLotusICEE_BoldIranic.ttf},
  FontFace={m}{sl}{IRLotusICEE_Iranic.ttf},
  Scale=1.15
]{IRLotusICEE.ttf}
\setlatintextfont[
  Path={../booklet/font/},
  BoldFont={LiberationSerif-Bold.ttf},
  BoldItalicFont={LiberationSerif-BoldItalic.ttf},
  ItalicFont={LiberationSerif-Italic.ttf}
]{LiberationSerif-Regular.ttf}
\setdigitfont[Path={../booklet/font/},Scale=1.15]{IRLotusICEE.ttf}
\defpersianfont\titlefont[Path={../booklet/font/},Scale=1]{IRTitr.ttf}

\definecolor{navy}{RGB}{20,65,105}
\definecolor{greenback}{RGB}{238,248,241}
\definecolor{yellowback}{RGB}{255,248,225}

\newcommand{\eng}[1]{\lr{#1}}
\newcommand{\term}[2]{\textbf{#1} \eng{(#2)}}
\newcommand{\key}[1]{\textcolor{navy}{\textbf{#1}}}

\newtcolorbox{formula}[1]{
  enhanced,breakable,colback=greenback,colframe=green!45!black,title={#1},
  fonttitle=\bfseries,boxrule=.6pt,arc=1.5mm,
  right=3.5mm,left=3.5mm,top=2mm,bottom=2mm,
  before skip=7pt,after skip=7pt,
  boxed title style={left=2mm,right=2mm,top=.7mm,bottom=.7mm}
}
\newtcolorbox{trap}[1]{
  enhanced,breakable,colback=yellowback,colframe=orange!65!black,title={#1},
  fonttitle=\bfseries,boxrule=.6pt,arc=1.5mm,
  right=3.5mm,left=3.5mm,top=2mm,bottom=2mm,
  before skip=7pt,after skip=7pt
}

\setlist[itemize,1]{
  label=\textcolor{navy}{\large\textbullet},
  rightmargin=7mm,leftmargin=2mm,topsep=4pt,itemsep=3.5pt,parsep=1pt
}
\setlist[itemize,2]{
  label=\textcolor{navy}{--},
  rightmargin=7mm,leftmargin=2mm,topsep=2pt,itemsep=2.5pt,parsep=0pt
}
\setlist[enumerate,1]{
  label=\textcolor{navy}{\arabic*.},
  rightmargin=9mm,leftmargin=2mm,topsep=4pt,itemsep=3pt
}
\setlength{\parindent}{0pt}
\setlength{\parskip}{3pt}
\setlength{\fboxsep}{5pt}
\linespread{1.12}
\renewcommand{\arraystretch}{1.3}

\titleformat{\section}[block]
  {\normalfont\titlefont\Large\color{white}}
  {}{0pt}
  {\colorbox{navy}{\parbox{\dimexpr\textwidth-2\fboxsep\relax}{\raggedleft\strut#1\strut}}}
\titlespacing*{\section}{0pt}{13pt}{10pt}

\titleformat{\subsection}[hang]
  {\normalfont\bfseries\large\color{navy}}
  {\thesubsection}{7pt}{#1}
  [\vspace{2pt}\color{navy!35}\titlerule]
\titlespacing*{\subsection}{0pt}{11pt}{7pt}

\pagestyle{fancy}
\fancyhf{}
\fancyhead[R]{\small مرور فصل ۸: هم‌ترازی مغز، کدگذاری عصبی و استدلال}
\fancyhead[L]{\small علوم شناختی محاسباتی}
\fancyfoot[C]{\thepage}
\setlength{\headheight}{14pt}

\begin{document}

\begin{center}
  {\titlefont\LARGE خلاصهٔ امتحانی فصل ۸}\\[3pt]
  {\Large هم‌ترازی مغز، کدگذاری عصبی و استدلال}\\[3pt]
  {\small مرور تفکیک‌شدهٔ نظریه‌ها، مراحل، مقایسه‌ها و فرمول‌ها}
\end{center}

\section{بخش اول: کدگذاری و رمزگشایی عصبی}

\subsection{\eng{Neural Encoding} و \eng{Neural Decoding}}

\begin{itemize}
  \item \term{کدگذاری عصبی}{Neural Encoding}: تبدیل محرک خارجی مانند صدا، تصویر یا لمس به فعالیت عصبی.
  \item \term{رمزگشایی عصبی}{Neural Decoding}: بازسازی محرک، نیت، فرمان، رفتار یا متن از روی فعالیت مغزی.
  \[
  \boxed{\text{\rl{محرک خارجی}}
  \longrightarrow\text{\rl{فعالیت مغزی}}
  \longrightarrow\text{\rl{فرمان یا رفتار}}}
  \]
  \item \eng{Encoding} مسیر اول و \eng{Decoding} تلاش برای پیمودن مسیر معکوس است.
\end{itemize}

\subsection{سه نظریهٔ اصلی \eng{Neural Coding}}

\begin{itemize}
  \item \term{\eng{Grandmother Cell}}{Grandmother Cell Theory}:
  \begin{itemize}
    \item یک نورون یا گروه بسیار کوچک، نمایندهٔ اختصاصی یک مفهوم است.
    \item مزیت: تفسیرپذیری بالا؛ ضعف: آسیب‌پذیری، ظرفیت نامعقول و نیاز به نماینده برای مفاهیم بسیار زیاد.
    \item شباهت محاسباتی: نودهای اختصاصی کلاس در لایهٔ خروجی.
  \end{itemize}
  \item \term{کدگذاری جمعیتی}{Population Coding}:
  \begin{itemize}
    \item تصمیم از مشارکت تعداد زیادی نورون با شدت و وزن متفاوت ساخته می‌شود:
    \[
    \boxed{\sum_i W_iX_i}
    \]
    \item مزایا: پایداری در برابر آسیب، دقت و ظرفیت ترکیبی بیشتر.
    \item ضعف‌ها: مصرف انرژی بیشتر، توزیع‌شدگی و تفسیرپذیری کمتر.
  \end{itemize}
  \item \term{کدگذاری تنک}{Sparse Coding}:
  \begin{itemize}
    \item فقط زیرمجموعه‌ای کوچک و مرتبط، مثلاً حدود \(5\%\) تا \(20\%\)، از نورون‌ها برای هر وظیفه فعال می‌شوند.
    \item حالت میانی: پایداری بهتر از بازنمایی اختصاصی، انرژی و تفسیرپذیری بهتر از کدگذاری کاملاً جمعیتی.
  \end{itemize}
\end{itemize}

\subsection{مقایسه و ترکیب نظریه‌ها}

\[
\boxed{
\begin{aligned}
\text{\rl{لایه‌های ابتدایی}}&\longrightarrow\mathrm{Population\ Coding}\\
\text{\rl{لایه‌های میانی}}&\longrightarrow\mathrm{Sparse\ Coding}\\
\text{\rl{لایه‌های انتهایی}}&\longrightarrow\mathrm{Grandmother\text{-}like}
\end{aligned}}
\]

\begin{itemize}
  \item ابتدایی: ویژگی‌های خام مانند فرکانس، واج، لبه، جهت و رنگ؛ توزیع‌شده و کم‌تفسیر.
  \item میانی: شکل و ویژگی‌های انتزاعی‌تر؛ گروه‌های منتخب فعال می‌شوند.
  \item انتهایی: کلاس یا مفهوم مشخص؛ بازنمایی اختصاصی‌تر.
  \item \eng{Symbolic AI} بیشتر به \eng{Sparse Coding} نزدیک است: مجموعه‌ای منتخب از قواعد و ویژگی‌های صریح در تصمیم نقش دارند.
\end{itemize}

\subsection{\eng{Brain-to-QWERTY} و \eng{Brain-to-Text}}

\begin{itemize}
  \item نمونهٔ مطرح‌شده: پژوهش \eng{Meta} در سال \(2025\) دربارهٔ \eng{Brain-to-QWERTY}.
  \item هدف محدود \eng{Brain-to-QWERTY}: رمزگشایی کلیدهایی که فرد هنگام تایپ فشار داده یا قصد فشار‌دادن آن‌ها را داشته است؛ نه خواندن تمام افکار.
  \item مرحلهٔ آموزش:
  \begin{enumerate}
    \item نمایش و خواندن متن توسط فرد.
    \item تایپ متن با صفحه‌کلید \eng{QWERTY}.
    \item ثبت هم‌زمان سیگنال مغزی.
    \item مقایسهٔ خروجی مدل با متن یا کلیدهای مورد انتظار.
    \item محاسبهٔ خطا، اصلاح وزن‌ها و کالیبراسیون برای همان فرد.
  \end{enumerate}
  \item فاز آزمون: متن هدف نامعلوم است و مدل فقط از فعالیت مغزی، متن یا نیت تایپ را بازسازی می‌کند.
\end{itemize}

\begin{formula}{معماری کلی \eng{Brain-to-Text}}
\[
\boxed{
\text{\rl{سیگنال مغزی}}
\rightarrow\mathrm{CNN}
\rightarrow\text{\rl{ویژگی میانی}}
\rightarrow\mathrm{Decoder\text{-}only\ Transformer}
\rightarrow\mathrm{Language\ Model}
\rightarrow\text{\rl{متن}}}
\]
\begin{itemize}
  \item \eng{CNN}: استخراج ویژگی از سیگنال خام.
  \item \eng{Transformer}: مدل‌کردن توالی کلیدها یا حروف.
  \item \eng{Language Model}: اصلاح کانتکستی و خطاهایی مانند «صلام» به «سلام».
\end{itemize}
\end{formula}

\subsection{کاربرد و مقایسه با \eng{Neuralink}}

\begin{itemize}
  \item کاربرد: ارتباط برای افراد دچار آسیب نخاعی یا محدودیت شدید حرکتی؛ تبدیل نیت ذهنی به متن یا فرمان.
  \item \textbf{\eng{Brain-to-QWERTY}}: دامنهٔ محدود به تایپ و کلیدهای صفحه‌کلید.
  \item \textbf{\eng{Neuralink}}: ایمپلنت تهاجمی برای تبدیل نیت حرکتی به کنترل ماوس، کلیک یا ابزار خارجی.
\end{itemize}

\subsection{روش‌های ثبت فعالیت مغزی}

\begin{itemize}
  \item \textbf{\eng{fMRI}:} غیرتهاجمی، مبتنی بر جریان خون، رزولوشن مکانی خوب و زمانی ضعیف در مقیاس ثانیه.
  \item \textbf{\eng{EEG}:} غیرتهاجمی، ثبت الکتریکی از پوست سر، رزولوشن زمانی خوب در حد میلی‌ثانیه و مکانی ضعیف‌تر.
  \item \textbf{ایمپلنت مغزی:} تهاجمی، نزدیک به نورون‌ها و دارای رزولوشن زمانی و مکانی بهتر؛ همراه با خطر جراحی، پایداری و مسائل اخلاقی.
\end{itemize}

\subsection{چالش‌ها و مسیر آینده}

\begin{itemize}
  \item \term{\eng{Plasticity}}{Plasticity}: مغز با یادگیری، خستگی و زمان تغییر می‌کند؛ مدل کالیبره‌شده ممکن است بعداً همان توزیع سیگنال را نبیند.
  \item \textbf{تعمیم بین افراد:} نقشهٔ مغزی مانند اثر انگشت متفاوت است؛ مدل یک فرد الزاماً روی دیگری کار نمی‌کند.
  \item \term{چندوجهی‌بودن}{Multimodality}: متن، تصویر، صدا، حافظه، توجه و هیجان هم‌زمان بر فعالیت مغز اثر دارند.
  \item هدف سخت‌افزاری: سنسور غیرتهاجمی، پوشیدنی، دقیق، کم‌مصرف و کم‌حرارت.
  \item \term{\eng{Neuromorphic}}{Neuromorphic Hardware}: پردازش رویدادمحور به‌جای پردازش دائم مبتنی بر کلاک:
  \[
  \boxed{\mathrm{Spike/Event}\longrightarrow\mathrm{Processing}}
  \]
\end{itemize}

\subsection{وزن نهایی نورونی}

\begin{formula}{عوامل شکل‌دهندهٔ وزن}
\[
\boxed{
\text{\rl{وزن نهایی}}
=\mathrm{Hebbian\ Learning}
+\text{\rl{ترجیح مکانی}}
+\text{\rl{مقداردهی اولیه}}}
\]
\begin{itemize}
  \item قانون \eng{Hebb}: هم‌فعالی مکرر، اتصال را تقویت می‌کند.
  \item ترجیح مکانی/عملکردی: نورون‌های سازگار با جهت یا کار خاص فعالیت بیشتری دارند.
  \item مقداردهی اولیه: تفاوت‌های اولیهٔ افراد بر پاسخ نهایی اثر می‌گذارد.
\end{itemize}
\end{formula}

\section{بخش دوم: \eng{Brain Alignment} و \eng{Brain Score}}

\subsection{ایدهٔ اصلی هم‌ترازی}

\begin{itemize}
  \item \term{فرضیهٔ بازنمایی افلاطونی}{Platonic Representation Hypothesis}: پشت بازنمایی‌های ظاهری، یک مفهوم یا حقیقت زیرین مشترک وجود دارد.
  \item این حقیقت زیرین را می‌توان با متغیر پنهان \(Z\) نشان داد؛ خود \(Z\) مستقیماً در دسترس نیست.
  \item تصویر، متن، صدا، رفتار، فعالیت مغزی و \eng{hidden state} مدل، بازنمایی‌های متفاوت همان \(Z\) مشترک‌اند:
  \[
  \boxed{
  Z\longrightarrow
  \left\{
  \text{\rl{متن، تصویر، رفتار، مغز، بازنمایی مدل}}
  \right\}}
  \]
  \item پاسخ درست به‌تنهایی فهم را ثابت نمی‌کند؛ مدل ممکن است الگوی سطحی را حفظ کرده باشد.
  \item مدل عمیق‌تر باید از پوستهٔ ظاهری عبور کند و ویژگی انتزاعی و رابطهٔ علّی را حفظ کند؛ مثلاً «سیب کپک‌زده» را با وجود تغییر ظاهر همچنان سیب تشخیص دهد.
  \item \textbf{\eng{Brain Alignment}}: سنجش اینکه بازنمایی‌های میانی مدل تا چه حد با بازنمایی‌های مغز انسان قابل نگاشت و هم‌راستا هستند.
  \item اثبات قطعی «فهم» ممکن نیست؛ هدف عملی، یافتن نقاطی است که مدل در آن‌ها از معنا یا رفتار انسان فاصله دارد.
\end{itemize}

\subsection{طراحی آزمایش هم‌ترازی}

\begin{enumerate}
  \item تعریف بخش‌های میانی مغز با \eng{Brain Parcellation}.
  \item انتخاب وظیفهٔ مناسب: زبانی، پیش‌بینی، استدلالی، علّی یا چندوجهی.
  \item ارائهٔ محرک یکسان به انسان و مدل.
  \item استخراج \eng{fMRI}/رفتار انسان و \eng{hidden state}های مدل.
  \item تعریف امتیاز با \eng{correlation}، \eng{similarity} یا \eng{loss}.
\end{enumerate}

\subsection{آزمایش \eng{Brain Score}}

\begin{itemize}
  \item مدل اصلی \eng{freeze} می‌شود؛ هدف ارزیابی بازنمایی موجود است.
  \item از لایه‌های مدل بردارهایی مانند \(v_1,v_2,v_3\) و از انسان \eng{fMRI} و رفتارهایی مانند زمان خواندن گرفته می‌شود.
  \item یک \term{مدل نگاشت}{Mapping Model} ساده، معمولاً خطی، یاد می‌گیرد:
  \[
  \boxed{\text{\rl{بازنمایی مدل}}\longrightarrow
  \text{\rl{سیگنال مغزی یا رفتار انسان}}}
  \]
  \item آموزش: یادگیری نگاشت روی متن‌های مشترک؛ آزمون: پیش‌بینی پاسخ مغز/رفتار برای متن جدید و مقایسه با دادهٔ واقعی.
  \item امتیاز می‌تواند ترکیبی از هم‌ترازی \eng{fMRI}، خروجی، زمان خواندن و پیش‌بینی باشد.
\end{itemize}

\subsection{اهمیت \eng{Next Word Prediction}}

\begin{itemize}
  \item وظایف بررسی‌شده: \eng{QA}، انسجام، \eng{entailment}، شباهت جمله و پیش‌بینی واژهٔ بعدی.
  \item همبستگی مثبت مهم‌تر در \eng{Next Word Prediction} دیده شد؛ مغز دائماً پیش‌بینی می‌کند.
  \item جملهٔ قابل‌پیش‌بینی: اطمینان مدل بیشتر و خواندن انسان سریع‌تر؛ عبارت غیرمنتظره: اطمینان کمتر و مکث بیشتر.
\end{itemize}

\subsection{\eng{Brain Parcellation}: واحدها}

\begin{itemize}
  \item \term{\eng{Voxel}}{Voxel}: واحد مکعبی اندازه‌گیری در \eng{fMRI}؛ نه یک مرز طبیعی زیستی.
  \item \term{\eng{Parcel}}{Parcel}: گروهی از وکسل‌ها با ویژگی مشترک.
  \item \term{\eng{Atlas}}{Atlas}: مجموعه‌ای از پارسل‌ها و نقشهٔ تقسیم‌بندی مغز.
  \item \term{\eng{ROI}}{Region of Interest}: ناحیهٔ منتخب متناسب با وظیفه، مثلاً نواحی زبان.
\end{itemize}

\[
\boxed{
\text{\rl{چند وکسل}}\longrightarrow
\text{\rl{یک پارسل}}\longrightarrow
\text{\rl{اطلس/افراز‌بندی}}}
\]

\subsection{سیگنال \eng{BOLD}}

\begin{itemize}
  \item \[
  \boxed{\mathrm{BOLD}=\mathrm{Blood\ Oxygen\ Level\ Dependent}}
  \]
  \item شاخص تغییر اکسیژن خون و نشانهٔ غیرمستقیم فعالیت عصبی.
  \item برای هر وکسل یک سری زمانی داریم؛ میانگین وکسل‌های یک پارسل، بازنمایی زمانی آن پارسل را می‌سازد.
  \item مثال: سه وکسل در چهار زمان \(\Rightarrow\) ماتریس \(3\times4\)، سپس بردار میانگین پارسل.
\end{itemize}

\subsection{روش‌های \eng{Parcellation}}

\begin{itemize}
  \item \textbf{آناتومیک:} تقسیم بر اساس محل و ساختار فیزیکی، مانند \eng{V1,V2,V4,IT}.
  \item \term{\eng{Functional Parcellation}}{Functional Parcellation}: گروه‌بندی نواحی دور یا نزدیک بر اساس هم‌فعال‌شدن و کارکرد مشترک.
  \item \term{\eng{Atlas}}{Atlas}: استفاده از تقسیم‌بندی ازپیش‌ساخته‌شده.
  \item \term{\eng{Localizer}}{Localizer}: یافتن ناحیه‌های فعال با اجرای همان وظیفه روی شرکت‌کننده.
\end{itemize}

\subsection{شبکه‌های هفت‌گانهٔ \eng{Yeo}}

\begin{enumerate}
  \item \textbf{\eng{Visual}:} پردازش دیداری، عمدتاً بخش پس‌سری.
  \item \textbf{\eng{Somatomotor}:} حس بدنی، حرکت و هماهنگی عضلات.
  \item \textbf{\eng{Dorsal Attention}:} توجه ارادی و درون‌هدایت‌شده.
  \item \textbf{\eng{Ventral Attention}:} توجه ناگهانی ناشی از محرک بیرونی.
  \item \textbf{\eng{Frontoparietal}:} کنترل، تصمیم‌گیری، هماهنگی چند وظیفه و رفتار.
  \item \textbf{\eng{Default Mode Network; DMN}:} استراحت، تفکر آزاد، خیال‌پردازی و ایده‌پردازی.
  \item \textbf{\eng{Limbic}:} احساس، هیجان و پردازش عاطفی.
\end{enumerate}

\subsection{\eng{Resting State} و همبستگی}

\begin{itemize}
  \item استراحت به معنای خاموشی مغز نیست؛ شبکه‌ها همچنان فعال و در حال پیش‌بینی‌اند.
  \item مغز حدود \(20\%\) انرژی بدن را مصرف می‌کند؛ وظیفهٔ شناختی بیشتر الگوی شبکه‌ها را تغییر می‌دهد و افزایش انرژی نسبت به خط پایه محدود است.
  \item ساخت افراز‌بندی کارکردی:
  \begin{enumerate}
    \item گرفتن سری زمانی \eng{BOLD} هر وکسل؛
    \item محاسبهٔ همبستگی سری‌ها؛
    \item خوشه‌بندی وکسل‌های با همبستگی بالا؛
    \item درنظرگرفتن هر خوشه به‌عنوان پارسل یا شبکه.
  \end{enumerate}
  \item \term{\eng{Task-Positive}}{Task-Positive}: هنگام انجام وظیفه فعال‌ترند؛ مثال‌ها: توجه پشتی، توجه شکمی و فرونتوپاریتال.
  \item \term{\eng{Task-Negative}}{Task-Negative}: با شروع وظیفه کاهش فعالیت دارند؛ مثال اصلی: \eng{DMN}.
\end{itemize}

\begin{trap}{مثال سریع \eng{Task-Positive} و \eng{Task-Negative}}
\begin{itemize}
  \item \textbf{مطالعهٔ متمرکز یا حل مسئله:} شبکه‌های توجه و فرونتوپاریتال بالا می‌روند و \eng{DMN} معمولاً کاهش می‌یابد.
  \item \textbf{استراحت، خیال‌پردازی و تفکر آزاد:} \eng{DMN} فعال‌تر است و شبکه‌های وظیفه‌مثبت فعالیت کمتری دارند.
  \item این رابطه مطلق و روشن/خاموشی نیست؛ الگوی فعالیت نسبی شبکه‌ها تغییر می‌کند.
\end{itemize}
\end{trap}

\subsection{بازگشت به \eng{Brain Alignment}}

\[
\boxed{
\mathrm{Parcel/BOLD}
\overset{\mathrm{Correlation/Similarity}}{\longleftrightarrow}
\mathrm{Hidden\ State}}
\]

\begin{itemize}
  \item همبستگی بالاتر یعنی هم‌ترازی بازنمایی بیشتر، نه اثبات قطعی فهم حقیقی.
  \item وظایف سخت‌تر مانند استنتاج، کشف قانون و تشخیص مغالطه برای آزمون معنا مناسب‌ترند.
\end{itemize}

\begin{enumerate}
  \item تقسیم مغز به پارسل‌های کارکردی.
  \item استخراج سیگنال \eng{BOLD} یا بازنمایی زمانی هر پارسل.
  \item استخراج \eng{hidden state}های لایه‌های میانی مدل.
  \item محاسبهٔ \eng{correlation} یا \eng{similarity} میان بازنمایی مغز و مدل.
\end{enumerate}

\section{بخش سوم: استدلال، قیاس و هم‌ترازی خطا}

\subsection{چرا پاسخ درست کافی نیست؟}

\begin{itemize}
  \item هم‌ترازی فقط شباهت پاسخ نهایی نیست؛ باید \key{الگوی خطا} و مسیر رسیدن به پاسخ نیز مقایسه شود.
  \item مدل هم‌تراز باید تا حدی در همان موقعیت‌هایی خطا کند که انسان‌ها خطا می‌کنند.
  \item دقت بیشتر از انسان الزاماً به معنای فهم یا استدلال انسانی نیست.
  \item برخی سوگیری‌ها و خطاهای شناختی را می‌توان \key{بهینه‌سازی تکاملی} دانست: فیلتر اطلاعات برای کاهش انرژی و تصمیم سریع‌تر برای بقا؛ بنابراین هر خطای انسانی صرفاً «باگ» نیست.
  \item \term{\eng{Pareidolia}}{Pareidolia}: دیدن چهره یا الگوی آشنا در ابر یا اشیای بی‌جان؛ خطای ادراکی که می‌تواند از نظر بقا و تعامل اجتماعی مفید باشد.
\end{itemize}

\subsection{\eng{System 1} و \eng{System 2}}

\begin{itemize}
  \item \textbf{\eng{System 1}:} سریع، کم‌هزینه، شهودی، مبتنی بر تجربه و \eng{prior knowledge}.
  \item \textbf{\eng{System 2}:} کند، پرهزینه، تحلیلی و نیازمند بررسی و اثبات.
  \item تفکر انتقادی یعنی در موقعیت لازم از پاسخ پیش‌فرض سیستم ۱ عبور و سیستم ۲ را فعال کنیم.
\end{itemize}

\subsection{انواع \eng{Reasoning}}

\begin{itemize}
  \item \term{قیاسی}{Deductive}: نتیجهٔ ضروری از مقدمات.
  \item \term{استقرایی}{Inductive}: تعمیم از نمونه‌ها؛ نتیجه احتمالی.
  \item \term{بهترین تبیین}{Abductive}: انتخاب محتمل‌ترین توضیح برای مشاهده.
  \item تمرکز فصل: \term{قیاس صوری}{Syllogistic Reasoning}.
\end{itemize}

\subsection{ساختار \eng{Syllogism}}

\begin{itemize}
  \item دو مقدمه و یک نتیجه:
  \[
  \mathrm{Premise\ 1}+\mathrm{Premise\ 2}
  \longrightarrow\mathrm{Conclusion}
  \]
  \item \term{\eng{Middle Term}}{Middle Term}: ترم مشترک میان دو مقدمه که آن‌ها را به هم متصل می‌کند؛ بدون آن قیاس شکل نمی‌گیرد.
\end{itemize}

\subsection{چهار \eng{Mode}}

\begin{itemize}
  \item \textbf{\eng{A} ـ کلی مثبت:} \(\mathrm{All\ A\ are\ B}\)
  \item \textbf{\eng{I} ـ جزئی مثبت:} \(\mathrm{Some\ A\ are\ B}\)
  \item \textbf{\eng{E} ـ کلی منفی:} \(\mathrm{No\ A\ are\ B}\)
  \item \textbf{\eng{O} ـ جزئی منفی:} \(\mathrm{Some\ A\ are\ not\ B}\)
\end{itemize}

\subsection{\eng{Figure}ها و تعداد فرم‌ها}

\[
\boxed{
\begin{array}{c@{\quad}c@{\quad}c@{\quad}c}
\mathrm{Figure\ 1}&\mathrm{Figure\ 2}&\mathrm{Figure\ 3}&\mathrm{Figure\ 4}\\
A-B&B-A&A-B&B-A\\
B-C&C-B&C-B&B-C
\end{array}}
\]

\begin{itemize}
  \item سه گزاره، هرکدام با چهار \eng{Mode} و چهار جایگاه ترم:
  \[
  \boxed{4\times4\times4=64\ \text{\rl{فرم ممکن}}}
  \]
  \item در صورت‌بندی فصل: \(27\) فرم دارای نتیجهٔ معتبر و \(37\) فرم دارای \term{NVC}{No Valid Conclusion}.
\end{itemize}

\subsection{نتیجهٔ معتبر و \eng{Barbara}}

\begin{itemize}
  \item \term{\eng{Valid Conclusion}}{Valid Conclusion}: نتیجه‌ای که در همهٔ حالت‌های ممکن درست باشد و هیچ مثال نقضی نداشته باشد.
\end{itemize}

\begin{formula}{فرم معتبر Barbara یا \((AAA)-1\)}
\[
\mathrm{All\ A\ are\ B}
\]
\[
\mathrm{All\ B\ are\ C}
\]
\[
\boxed{\therefore\ \mathrm{All\ A\ are\ C}}
\]
\end{formula}

\begin{itemize}
  \item درست‌بودن نتیجه در بعضی حالت‌ها کافی نیست؛ باید در \key{تمام} حالت‌ها برقرار باشد.
  \item نمودار ون ابزار بررسی وجود مثال نقض و اعتبار نتیجه است.
\end{itemize}

\subsection{قیاس به‌عنوان آزمون \eng{Brain Alignment}}

\begin{itemize}
  \item هر ۶۴ فرم به انسان و مدل داده می‌شود؛ سپس دقت و مهم‌تر از آن، توزیع پاسخ و خطا مقایسه می‌شود.
  \item اگر مدل همیشه درست باشد اما در فرم‌های دشوار شبیه انسان خطا نکند، الزاماً هم‌تراز نیست.
  \item انسان برای تمام فرم‌های منطقی بهینه نشده؛ زندگی واقعی بیشتر به تصمیم سریع، تجربه، بایاس و مغالطه وابسته است.
\end{itemize}

\subsection{زبان طبیعی و دیتاست \eng{Avicenna}}

\begin{itemize}
  \item تبدیل \(A,B,C\) به زبان طبیعی مسئله را دشوار می‌کند: مدل باید ترم میانی، رابطهٔ منطقی، مغالطه و \eng{NVC} را تشخیص دهد.
  \item دیتاست \eng{Avicenna}: حدود \(6000\) نمونهٔ قیاس در زبان طبیعی.
  \item اجزای هر نمونه:
  \begin{itemize}
    \item \eng{Premise 1} و \eng{Premise 2}
    \item \eng{Syllogistic Relation}
    \item \eng{Conclusion} یا \eng{NVC}
  \end{itemize}
  \item اشتراک واژگانی الزاماً ترم میانی نیست؛ «شیر جنگل» و «شیر آب» واژهٔ مشترک اما معنای متفاوت دارند.
\end{itemize}

\subsection{\eng{Syllogism} در برابر \eng{Entailment}}

\begin{itemize}
  \item \term{\eng{Syllogism}}{Syllogism}: دو مقدمه با ترم میانی مشترک؛ نتیجه از ترکیب آن‌ها حاصل می‌شود.
  \item \term{\eng{Entailment}}{Entailment}: ممکن است یک گزاره و دانش عمومی برای نتیجه‌گیری کافی باشد.
  \item هر استنتاج زبانی الزاماً قیاس نیست.
\end{itemize}

\subsection{\eng{Prior Knowledge} و مغالطه}

\begin{itemize}
  \item «دولت فاسد تورم ایجاد می‌کند؛ ایران تورم دارد؛ پس دولت ایران فاسد است» نامعتبر و شبیه \term{تأیید تالی}{Affirming the Consequent} است.
  \item نارضایتی یا باور قبلی می‌تواند نتیجهٔ نامعتبر را برای سیستم ۱ پذیرفتنی کند.
  \item برعکس، یک نتیجهٔ قیاسی معتبر دربارهٔ رشد اقتصادی ممکن است به‌علت تجربهٔ شخصی پذیرفته نشود.
  \item بنابراین اعتبار منطقی و پذیرش روان‌شناختی یکسان نیستند.
\end{itemize}

\subsection{مدل در برابر انسان}

\begin{itemize}
  \item مدل بزرگ‌تر ممکن است در فرم‌های صوری از انسان دقیق‌تر باشد، اما توزیع خطایش متفاوت بماند.
  \item پاسخ بهتر می‌تواند حاصل دادهٔ بیشتر، حافظه یا دیدن نمونه‌های مشابه باشد؛ نه الزاماً فهم انسانی.
  \item انسان و مدل هر دو آماری یاد می‌گیرند، اما توزیع تجربه و خطاهای آن‌ها یکسان نیست.
  \item \term{\eng{Adversarial Data Augmentation}}{Adversarial Data Augmentation}: شناسایی فرم‌های شکست مدل و افزودن هدفمند نمونه‌های مشابه؛ روشی استفاده‌شده برای بهبود روی \eng{Avicenna}.
\end{itemize}

\begin{trap}{معیار نهایی هم‌ترازی}
برای سنجش مدل فقط دقت کافی نیست. باید پاسخ درست، خطاهای انسان، خطاهای مدل، توزیع خطا، بازنمایی‌های میانی، اثر دانش قبلی و تفاوت استدلال صوری با استدلال جهان واقعی را هم‌زمان بررسی کرد.
\end{trap}

\end{document}
